% ===== Block B07: 原文 p.54–66 | §3.3 共形块空间的基本性质 | 译手 tr-ch3a =====
\subsection{共形块空间的基本性质}

% [原文 p.54]
本节我们将说明共形块空间的基本性质。

\noindent\textbf{3.3.1 真空传播.}\quad
首先我们证明所谓的\emph{真空传播}，它在描述真空空间时起重要作用。

对 $\mathfrak{X}=(C;Q_1,\dots,Q_N;\eta_1,\dots,\eta_N)$，设 $P$ 是曲线 $C$ 的非奇异点，$t$ 是 $C$ 在 $P$ 处的形式参数。令
\[
\tilde{\mathfrak{X}}=(C;Q_1,\dots,Q_N,Q_{N+1};\eta_1,\dots,\eta_N,\eta_{N+1}),
\]
其中 $Q_{N+1}=P$，$\eta_{N+1}=\eta$。

以下我们固定可积左 $\hat{\mathfrak{g}}$-模 $\cH_0$ 的一个最高权向量 $|0\rangle$。由于存在典范包含
\[
\begin{aligned}
\cH_{\mathfrak{X}}&\longrightarrow\cH_{\mathfrak{X}}\otimes\cH_0,\\
|v\rangle&\longmapsto|v\rangle\otimes|0\rangle,
\end{aligned}
\]
我们有典范满射
\[
\hat\pi^*:\cH^{\dagger}_{\mathfrak{X}}\widehat{\otimes}\cH^{\dagger}_0\longrightarrow\cH^{\dagger}_{\mathfrak{X}}.
\]

\begin{theorem}[3.15]
上述典范满射 $\hat\pi^*$ 诱导出典范同构
\[
\cV^{\dagger}_{\vec{\lambda},0}(\tilde{\mathfrak{X}})\cong\cV^{\dagger}_{\vec{\lambda}}(\mathfrak{X}).
\]
\end{theorem}

\noindent\textit{证明.}\quad 对元素 $\langle\tilde\Psi|\in\cV^{\dagger}_{\vec{\lambda},0}(\tilde{\mathfrak{X}})$，令 $\langle\Psi|=\hat\pi^*(\langle\tilde\Psi|)\in\cH^{\dagger}_{\mathfrak{X}}$。取 $f\in H^0(C,\cO_C(*\sum_{j=1}^{N}Q_j))$、$X\in\mathfrak{g}$ 以及 $|u\rangle\in\cH_{\mathfrak{X}}$。则由我们的定义有
\[
\sum_{j=1}^{N}\langle\Psi|\rho_j(X[f])|u\rangle=\sum_{j=1}^{N}\langle\tilde\Psi|\rho_j(X[f])|u\otimes0\rangle.
\]
另一方面，由于 $f$ 在点 $Q_{N+1}=P$ 处是正则的，我们有
\[
\langle\tilde\Psi|\rho_{N+1}(X[f])|u\otimes0\rangle=0.
\]

% [原文 p.55]
\figomit{3.1}{真空传播。}

因此我们有
\[
\sum_{j=1}^{N}\langle\tilde\Psi|\rho_j(X[f])|u\otimes0\rangle
=\sum_{j=1}^{N+1}\langle\tilde\Psi|\rho_j(X[f])|u\otimes0\rangle=0.
\]
于是 $\langle\Psi|\in\cV^{\dagger}_{\vec{\lambda}}(\mathfrak{X})$，并且我们得到线性映射
\[
\iota^*:\cV^{\dagger}_{\vec{\lambda},0}(\tilde{\mathfrak{X}})\longrightarrow\cV^{\dagger}_{\vec{\lambda}}(\mathfrak{X}).
\]
首先我们证明线性映射 $\iota^*$ 是单射。

假设 $\langle\Psi|=\iota^*(\langle\tilde\Psi|)=0$。通过对 $p$ 的归纳，我们证明
\begin{equation}
\langle\tilde\Psi|u\otimes v\rangle=0,\qquad\text{对一切 }u\in\cH_{\mathfrak{X}}\text{ 与 }v\in F_p\cH_0. \tag{3.24}
\end{equation}
由我们的假设有
\[
\langle\tilde\Psi|u\otimes0\rangle=\langle\Psi|u\rangle=0.
\]
因此 (3.24) 对 $p=0$ 成立。其次假设 (3.24) 对 $p$ 成立。取元素 $X(m)|v\rangle\in F_{p+1}\cH_0$，其中 $|v\rangle\in F_p\cH_0$。取亚纯函数 $f\in H^0(C,\cO_C(*\sum_{j=1}^{N+1}Q_j))$ 与正整数 $M$，使得在点 $P$ 处有
\begin{equation}
f\equiv\eta^{m}\mod(\eta^{M}) \tag{3.25}
\end{equation}
并且使得
\begin{equation}
X\otimes\eta^{k}|v\rangle=0\qquad\text{对一切 }k\ge M. \tag{3.26}
\end{equation}
于是我们有
\[
\begin{aligned}
\langle\tilde\Psi|u\otimes X(m)|v\rangle
&=\langle\tilde\Psi|u\otimes(X[f])v\rangle\\
&=-\sum_{j=1}^{N}\langle\tilde\Psi|\rho_j(X[f])u\otimes v\rangle\\
&=0,
\end{aligned}
\]
因为由归纳假设 $\sum_{j=1}^{N}\langle\tilde\Psi|\rho_j(X[f])u\otimes v\rangle=0$。因此 (3.24) 对 $p+1$ 成立。于是对任意 $|u\otimes v\rangle\in\cH_{\mathfrak{X}}\otimes\cH_0$ 都有 $\langle\tilde\Psi|u\otimes v\rangle=0$。因此 $\langle\tilde\Psi|=0$。

\begin{sloppypar}下面我们证明 $\iota^*$ 是满射。为此，对给定的 $\langle\Psi|\in\cV^{\dagger}_{\vec{\lambda}}(\mathfrak{X})$，我们对应一个元素 $\langle\tilde\Psi|\in\Hom_{\C}(\cH_{\mathfrak{X}}\otimes\cM_0,\C)$，其中 $\cM_0$ 是 Verma
\end{sloppypar}

% [Tutor Note: 单射性证明的关键：$\langle\tilde\Psi|$ 在 $\cV^\dagger$ 中意味着它对 $\hat{\mathfrak{g}}(\tilde{\mathfrak{X}})$ 的作用为零（(3.11)），而 $f\in H^0(C,\cO_C(*\sum_{j=1}^{N+1}Q_j))$ 在 $P$ 处有极点，使 $X[f]$ 把 $Q_{N+1}$ 处的分量也带上；于是
%   $\langle\tilde\Psi|u\otimes(X[f])v\rangle$ 可用「$\hat{\mathfrak{g}}(\tilde{\mathfrak{X}})$-不变性」写成 $-\sum_{j=1}^{N}\langle\tilde\Psi|\rho_j(X[f])u\otimes v\rangle$，再对 $p$ 归纳。】
%   跳步在于：证明 $\langle\tilde\Psi|=0$ 需要它在 $\cH_{\mathfrak{X}}\otimes\cH_0$ 上为零，而这里只对 $v\in F_p\cH_0$ 归纳；由于 $\cH_0$ 由 $|0\rangle$ 生成且滤过穷竭，归纳覆盖全空间。]

% [原文 p.56]
模，它是与 $\mathfrak{g}$ 的平凡表示相伴的仿射李代数 $\hat{\mathfrak{g}}$ 的模。（见 (2.4)。）

线性泛函 $\langle\tilde\Psi|$ 如下归纳地定义为 $\cH_{\mathfrak{X}}\otimes F_p\cM_0$ 到 $\C$ 的线性映射。首先定义
\[
\langle\tilde\Psi|u\otimes0\rangle=\langle\Psi|u\rangle
\qquad\text{对任意 }u\in\cH_{\mathfrak{X}}.
\]
于是我们有
\[
\sum_{j=1}^{N}\langle\tilde\Psi|\rho_j(X[g])|u\otimes0\rangle
=\sum_{j=1}^{N}\langle\Psi|\rho_j(X[g])|u\rangle=0
\]
对任意元素 $g\in H^0(C,\cO_C(*\sum_{j=1}^{N}Q_j))$ 成立。

现在假设 $\langle\tilde\Psi|$ 已定义为 $\cH_{\mathfrak{X}}\otimes F_p\cM_0$ 到 $\C$ 的线性映射，并且满足
\begin{equation}
\sum_{j=1}^{N}\langle\tilde\Psi|\rho_j(X[g])|u\otimes v\rangle=0 \tag{3.27}
\end{equation}
对任意 $|u\otimes v\rangle\in\cH_{\mathfrak{X}}\otimes F_p\cM_0$ 与 $g\in H^0(C,\cO_C(*\sum_{j=1}^{N}Q_j))$ 成立。那么在 $\cH_{\mathfrak{X}}\otimes F_{p+1}\cM_0$ 上，线性映射 $\langle\tilde\Psi|$ 由
\begin{equation}
\langle\tilde\Psi|u\otimes X(m)v\rangle
=-\sum_{j=1}^{N}\langle\tilde\Psi|(\rho_j(X[f])u\otimes v\rangle
\qquad\text{对任意 }u\in\cH_{\mathfrak{X}} \tag{3.28}
\end{equation}
定义，其中 $X(m)|v\rangle\in F_{p+1}\cM_0$、$|v\rangle\in F_p\cM_0$，而亚纯函数 $f$ 按 (3.25) 与 (3.26) 中同样的方式选取。容易证明这是良定义的，并且具有性质
\[
\sum_{j=1}^{N+1}\langle\tilde\Psi|\rho_j(X[g])=0
\qquad\text{在 }\cH_{\mathfrak{X}}\otimes F_{p+1}\cM_0\text{ 上}
\]
对每个元素 $g\in H^0(C,\cO_C(*\sum_{j=1}^{N+1}Q_j))$ 成立。一个直接的计算给出等式
\[
\begin{aligned}
&\langle\tilde\Psi|u\otimes X(m_1)Y(m_2)v\rangle
-\langle\tilde\Psi|u\otimes Y(m_2)X(m_1)v\rangle\\
&\qquad=\langle\tilde\Psi|u\otimes\bigl([X,Y](m_1+m_2)
+\ell\cdot(X,Y)m_1\delta_{m_1+m_2,0}\bigr)v\rangle
\end{aligned}
\]
只要 $X(m_1)Y(m_2)|v\rangle\in F_{p+1}\cM_0$ 且 $|v\rangle\in F_{p+1}\cM_0$。这个等式表明 $\langle\tilde\Psi|$ 已定义为 $\cH_{\mathfrak{X}}\otimes\cM_0$ 到 $\C$ 的线性映射。为证明 $\langle\tilde\Psi|$ 能下降为 $\cH_{\mathfrak{X}}\otimes\cH_0$ 上的线性形式，只需证明等式
\begin{equation}
\langle\tilde\Psi|u\otimes X_\theta(-1)^{\ell+1}|0\rangle=0. \tag{3.29}
\end{equation}
为证明 (3.29)，我们首先证明
\[
\langle\tilde\Psi|u\otimes X_\theta(-1)^{n}|0\rangle=0
\]
对充分大的 $n$（依赖于 $|u\rangle$）成立。设 $f\in H^0(C,\cO_C(*\sum_{j=1}^{N+1}Q_j))$ 是 $C$ 上的亚纯函数，满足 $m=-1$ 时的条件 (3.25) 与 (3.26)。由定理 3.15，存在依赖于 $|u\rangle$ 的正整数 $n$，使得对任意 $j$，$j=1,\dots,N$，都有
\begin{equation}
\rho_j(X_\theta[f])^{k}|u\rangle=0,\qquad\text{若 }k\ge n/N. \tag{3.30}
\end{equation}

% [Tutor Note: 这里 $\cM_0$ 是 $\hat{\mathfrak{g}}$ 对应于 $\mathfrak{g}$ 的平凡表示（即 $0$ 权）的 Verma 模，$\cH_0$ 是它的不可约商；先在线性泛函层面把 $\langle\tilde\Psi|$ 造在 $\cH_{\mathfrak{X}}\otimes\cM_0$ 上，再验证它穿过 $\cH_0$，是本节所有证明的通用策略。]
% [Tutor Note: (3.28) 的良定义性靠的是：若 $X(m)|v\rangle=0$，右端也应为 $0$；这由 $f$ 的选取 (3.25)(3.26) 与 (3.27) 保证。读者若跳过「well-defined」一句，会看不出为何 $f$ 要满足 $f\equiv\eta^m\bmod(\eta^M)$。]
% [Tutor Note: 原文「By Theorem 3.15 there is a positive integer $n$…」中的定理编号可疑：定理 3.15 是真空传播（$\cV^\dagger_{\vec\lambda,0}(\tilde{\mathfrak{X}})\cong\cV^\dagger_{\vec\lambda}(\mathfrak{X})$），给不出 $\rho_j(X_\theta[f])$ 的幂零性；此处逻辑上需要的应是定理 2.4（可积模的局部幂零性）。照原文排印，待核。]

% [原文 p.57]
应用公式 (3.28)，由 (3.30) 我们得到
\begin{equation}
\begin{aligned}
\langle\tilde\Psi|u\otimes X_\theta(-1)^{n}|0\rangle
&=\langle\tilde\Psi|u\otimes(X_\theta[f])^{n}|0\rangle\\
&=(-1)^{n}\sum_{n_1+\cdots+n_N=n}
\frac{n!}{n_1!n_2!\cdots n_N!}
\langle\tilde\Psi|\prod_{j=1}^{N}\rho_j(X_\theta[f])^{n_j}u\otimes0\rangle\\
&=0.
\end{aligned}\tag{3.31}
\end{equation}
令
\[
E=X_{-\theta}(1),\qquad F=X_\theta(-1),\qquad H=[E,F].
\]
则 $\{E,F,H\}$ 张成 $\hat{\mathfrak{g}}$ 的一个李子代数 $\mathfrak{k}$，同构于 $\mathfrak{sl}(2,\C)$。注意我们有
\[
H=[X_{-\theta},X_\theta]\otimes1+(X_{-\theta},X_\theta)c=-H_\theta\otimes1+c.
\]
因此有
\[
H|0\rangle=\ell|0\rangle.
\]
设 $U$ 是 Verma 模 $\cM_0$ 的由
\[
U:=\sum_{k=1}^{\infty}\C X_\theta(-1)^{k}|0\rangle
\]
定义的向量子空间。

通过简单的计算我们有
\[
HF^{k}|0\rangle=(\ell-2k)F^{k}|0\rangle,\qquad
EF^{k}|0\rangle=(k\ell-k^{2}+k)F^{k-1}|0\rangle.
\]
因此，向量空间 $U$ 是一个 $\mathfrak{k}$-模。$U$ 的子空间 $U_0$ 定义为满足
\[
\langle\tilde\Psi|u\otimes v\rangle=0,
\]
对一切 $|u\rangle\in\cH_{\mathfrak{X}}$ 成立的向量 $|v\rangle\in U$ 的集合。于是由 (3.31)，$U_0$ 是一个 $\mathfrak{k}$-模。由于对充分大的 $m$，由 (3.30) 我们有 $F^{m}|0\rangle\in U_0$，所以商空间 $U/U_0$ 是有限维 $\mathfrak{k}$-模。设 $|\overline{0}\rangle$ 是 $|0\rangle\in U_0$ 在 $U/U_0$ 中的像。于是我们有
\[
H|\overline{0}\rangle=\ell|\overline{0}\rangle.
\]
因此，由 $\mathfrak{sl}(2,\C)$ 的表示论我们有
\[
F^{\ell+1}|\overline{0}\rangle=0.
\]
这意味着
\[
\langle\tilde\Psi|u\otimes X_\theta(-1)^{\ell+1}|0\rangle=0.
\]
于是我们得到 $\langle\tilde\Psi|\in\cV^{\dagger}_{\vec{\lambda},0}(\tilde{\mathfrak{X}})$，使得 $\iota^*(\langle\tilde\Psi|)=\langle\Psi|$。\hfill$\square$

\begin{corollary}[3.16]
存在典范同构
\[
\cV_{\vec{\lambda}}(\mathfrak{X})\cong\cV_{\vec{\lambda},0}(\tilde{\mathfrak{X}}).
\]
\end{corollary}

\begin{remark}[3.17]
定理 3.15 与推论 3.16 表明，在研究带形式邻域的 $N$ 点稳定曲线所附着的共形块与共真空空间时，我们可以添加任意多个带形式邻域的点。因此，正如我们上面提到的，\emph{我们总可以假设条件 (Q) 成立}。下文将经常使用这一事实，它在证明若干重要结论时将起本质作用。
\end{remark}

在证明上述定理时，我们本质上用到了下面的定理。

% [Tutor Note: (3.31) 的推导要点：把 $X_\theta(-1)^n|0\rangle$ 换成 $(X_\theta[f])^n|0\rangle$（在 $P$ 处 $f\equiv\eta^{-1}$，故二者在 $P$ 处相差正则项，作用在 $|0\rangle$ 上相同），再用 (3.28) 逐次把 $\hat{\mathfrak{g}}(\tilde{\mathfrak{X}})$ 的作用移到左边，$n$ 次后得到 $(-1)^n$ 与多重和；由于和式中每个 $n_j$ 都可能很大，而 (3.30) 只保证 $\rho_j(X_\theta[f])^{k}|u\rangle=0$ 当 $k\ge n/N$，故至少有一个 $n_j$ 满足条件，从而整个和为零。
%   原文「the image of $|0\rangle\in U_0$ to $U/U_0$」中 $|0\rangle\in U_0$ 与此处语境不符：若 $|0\rangle\in U_0$，则它在 $U/U_0$ 中的像为 $0$，后面的 $H|\overline{0}\rangle=\ell|\overline{0}\rangle$ 便不成立（除非 $\ell=0$）。按上下文应为 $|0\rangle\in U$。照原文排印，登记为存疑。]
% [Tutor Note: 这里用到 $\mathfrak{sl}(2,\C)$ 的有限维表示论：$H$ 的特征值为 $\ell,\ell-2,\dots$，若 $F^{\ell+1}|\overline{0}\rangle\ne0$，则 $|\overline{0}\rangle,F|\overline{0}\rangle,\dots,F^{\ell}|\overline{0}\rangle$ 线性无关（权重各不相同），给出 $\ell+1$ 维不可约表示，而 $\ell$ 不是非负整数时会与有限维性矛盾；故 $F^{\ell+1}|\overline{0}\rangle=0$，再回到 (3.29)。]

% [原文 p.58]
\begin{theorem}[3.18]
典范包含
\[
V_{\lambda_1}\otimes\cdots\otimes V_{\lambda_{N-1}}\otimes\cH_{\lambda_N}
\hookrightarrow\cH_{\vec{\lambda}}
\]
诱导出自然投影
\[
p:\cV^{\dagger}_{\vec{\lambda}}\longrightarrow
\Hom\bigl(V_{\lambda_1}\otimes\cdots\otimes V_{\lambda_{N-1}}
\otimes\cH_{\lambda_N},\C\bigr)
=\cV^{\dagger}_{\lambda_1}\otimes\cdots\otimes
\cV^{\dagger}_{\lambda_{N-1}}\otimes\cH^{\dagger}_{\lambda_N}.
\]
其像包含于
\[
V^{\dagger,0}_{\vec{\lambda}}
=\bigl\{\langle\Phi|\in\cV^{\dagger}_{\lambda_1}\otimes\cdots\otimes
\cV^{\dagger}_{\lambda_{N-1}}\otimes\cH^{\dagger}_{\lambda_N}
\mid\langle\Phi|X\otimes f=0,\quad
\forall f\in H^0(C,\cO_C(*Q_N))\bigr\},
\]
并且映射 $p:\cV^{\dagger}_{\vec{\lambda}}\to V^{\dagger,0}_{\vec{\lambda}}$
是同构。
\end{theorem}

\noindent\textit{证明.}\quad 只需证明
$\langle\Phi|\in V^{\dagger,0}_{\vec{\lambda}}$ 可以唯一地延拓为
$\cV^{\dagger}_{\vec{\lambda}}$ 的元素。当 $N=1$ 时结论平凡。设
$N\ge2$。首先我们证明
$\langle\Phi|\in V^{\dagger,0}_{\vec{\lambda}}$ 可以延拓为
\[
\Hom\bigl(V_{\lambda_1}\otimes\cdots\otimes V_{\lambda_{N-2}}
\otimes\cM\otimes\cH_{\lambda_N},\C\bigr)
\]
的唯一元素 $\langle\tilde\Phi|$，使得
\[
\langle\tilde\Phi|X[f]=0
\qquad\text{对一切 }f\in H^0(C,\cO_C(*Q_{N-1}+Q_N)),
\]
其中 $\cM$ 是 Verma 模 $\cM_{\lambda_{N-1}}$。假设 $\langle\Phi|$ 已延拓为唯一元素
\[
\langle\Phi^{(p)}|\in\Hom\bigl(V_{\lambda_1}\otimes\cdots\otimes
V_{\lambda_{N-2}}\otimes F_p\cM\otimes\cH_{\lambda_N},\C\bigr),
\]
使得
\[
\langle\Phi^{(p)}|X[f]|\psi\rangle=0,
\]
只要 $f\in H^0(C,\cO_C(pQ_{N-1}+*Q_N))$，而 $|\psi\rangle$ 与
$X[f]|\psi\rangle$ 都属于
$V_{\lambda_1}\otimes\cdots\otimes V_{\lambda_{N-2}}\otimes F_p\cM
\otimes\cH_{\lambda_N}$。对 $p=0$，置
$\langle\Phi^{(0)}|=\langle\Phi|$。现在假设
$\langle\Phi^{(p)}|$ 已定义。

取元素 $|X(-n)\phi\rangle\in F_{p+1}\cM$，$n\ge1$，其中
$|\phi\rangle\in F_p\cM$。设 $m$ 是正整数，使得
\[
X(k)|\phi\rangle=0,\qquad\text{对一切 }k\ge m.
\]
然后选取亚纯函数
$g\in H^0(C,\cO_C(nQ_{N-1}+*Q_N))$，使得它在
$Q_{N-1}$ 处关于 $\xi_{N-1}$ 的 Laurent 展开具有形式
\begin{equation}
g_{N-1}=\frac{1}{\xi_{N-1}^{n}}+a_m\xi_{N-1}^{m}+\cdots. \tag{3.32}
\end{equation}
现在对元素
$|v_1\otimes\cdots\otimes v_{N-2}\otimes X(-n)\phi\otimes\psi\rangle$
定义
\[
\begin{aligned}
&\langle\Phi^{(p+1)}|v_1\otimes\cdots\otimes v_{N-2}
\otimes X(-n)\phi\otimes\psi\rangle\\
&\qquad=\sum_{i\ne N-1}\langle\Phi^{(p)}|
\rho_i(X[g])|v_1\otimes\cdots\otimes v_{N-2}
\otimes X(-n)\phi\otimes\psi\rangle.
\end{aligned}
\]
右端是良定义的，且与满足 (3.32) 的函数
$g\in H^0(C,\cO_C(nQ_{N-1}+*Q_N))$ 的选取无关。于是我们有
\[
\langle\Phi^{(p+1)}|X[g]|v_1\otimes\cdots\otimes v_{N-2}
\otimes\phi\otimes\psi\rangle=0.
\]
由于 $F_{p+1}\cM$ 的任一元素都是 $X_j(-n_j)\phi_j$（$j=1,\dots$，其中
$X_j\in\mathfrak{g}$）的有限线性组合，$\langle\Phi^{(p+1)}|$
因此被唯一确定。

让我们选取 $f\in H^0(C,\cO_C((p+1)Q_{N-1}+*Q_N))$ 与
$|\phi_1\rangle\in F_{p+1}\cM$，使得
$\rho_{N-1}(X[f])|\phi_1\rangle\in F_{p+1}\cM$。设
\[
\frac{b_{-p-1}}{\xi_{N-1}^{p+1}}+\frac{b_{-p}}{\xi_{N-1}^{p}}+\cdots
\]

% [Tutor Note: 定理 3.18 是本节的「投影」机制，务必与 (3.9)(3.10) 对照理解：$\cV^\dagger_{\vec\lambda}$ 是 $\cH^\dagger_{\vec\lambda}$ 中同时被 $\hat{\mathfrak{g}}(\mathfrak{X})$ 作用零化的元素，而典范包含 $V_{\lambda_1}\otimes\cdots\otimes\cH_{\lambda_N}\hookrightarrow\cH_{\vec\lambda}$ 的对偶给出投影 $p$。$p$ 的像是那些「只对最后一个因子上的 $\hat{\mathfrak{g}}(\mathfrak{X})$-作用零化」的泛函；定理断言这两步之间没有信息损失（$p$ 是同构）。]
% [Tutor Note: 证明是「从 $F_p$ 到 $F_{p+1}$ 逐层延拓」：先在 $F_0$ 上令 $\langle\Phi^{(0)}|=\langle\Phi|$，再用 $g$ 在 $Q_{N-1}$ 处的主部 $1/\xi_{N-1}^{n}$ 把 $X(-n)\phi$ 上的值化为 $\langle\Phi^{(p)}|$ 在其它点上的值；右端与 $g$ 的选取无关这一点，靠的是两个这样的 $g$ 之差在 $Q_{N-1}$ 处正则，从而由 $\langle\Phi^{(p)}|X[f]|\psi\rangle=0$ 消掉。]
% [Tutor Note: (3.32) 中只写出主部 $1/\xi_{N-1}^{n}$ 与首项 $a_m\xi_{N-1}^{m}$，中间各项被省略；$m$ 的选取使 $X(k)|\phi\rangle=0$（$k\ge m$），故 $\xi_{N-1}^{m}$ 及更高次项作用在 $|\phi\rangle$ 上为零，$\langle\Phi^{(p)}|$ 只需用到这些项。]

% [原文 p.59]
为 $f$ 在 $Q_{N-1}$ 处关于 $\xi_{N-1}$ 的 Laurent 展开。我们可以假设
$b_{-p-1}\ne0$。于是 $|\phi_1\rangle$ 属于 $F_p\cM$。现在选取正整数
$m_1$，使得对任意 $k\ge m_1$ 都有 $X(k)|\phi_1\rangle=0$；并选取
$h\in H^0(C,\cO_C((p+1)Q_{N-1}+*Q_N))$，使得 $h$ 在 $Q_{N-1}$ 处关于
$\xi_{N-1}$ 的 Laurent 展开具有形式
\[
\frac{1}{\xi_{N-1}^{p+1}}+c_{m_1}\xi_{N-1}^{m_1}+\cdots.
\]
于是由 $\langle\Phi^{(p+1)}|$ 的上述定义我们有
\[
\langle\Phi^{(p+1)}|X[h]v_1\otimes\cdots\otimes v_{N-2}
\otimes\phi_1\otimes\psi\rangle=0.
\]
此外我们有 $f-b_{-p-1}h\in H^0(C,\cO_C((p)Q_{N-1}+*Q_N))$，并且
\[
\rho_{N-1}(X[f-b_{-p-1}h])|\phi_1\rangle\in F_p\cM.
\]
于是由归纳假设我们有
\[
\langle\Phi^{(p+1)}|X[f-b_{-p-1}h]|v_1\otimes\cdots\otimes v_{N-2}
\otimes\phi_1\otimes\psi\rangle=0.
\]
这意味着
\[
\langle\Phi^{(p+1)}|X[f]v_1\otimes\cdots\otimes v_{N-2}
\otimes\phi\otimes\psi\rangle=0.
\]
于是我们已经证明 $\langle\Phi|$ 可以唯一地延拓为 $\langle\tilde\Phi|$。

下面我们证明 $\langle\tilde\Phi|$ 下降为
\[
\Hom\bigl(V_{\lambda_1}\otimes\cdots\otimes V_{\lambda_{N-2}}
\otimes\cH_{\lambda_{N-1}}\otimes\cH_{\lambda_N},\C\bigr)
\]
的元素。为此只需证明
\[
\langle\tilde\Phi|v_1\otimes\cdots\otimes v_{N-2}
\otimes X_\theta(-1)^{\ell-(\theta,\lambda_{N-1})+1}\lambda_{N-1}
\otimes\psi\rangle=0
\]
对任意 $v_i\in V_i$、$i=1,\dots,N-1$、$|\psi\rangle\in\cH_{\lambda_N}$
以及最高权向量 $\lambda_{N-1}$ 成立。首先我们证明
\[
\langle\tilde\Phi|v_1\otimes\cdots\otimes v_{N-2}
\otimes X_\theta(-1)^{M}\lambda_{N-1}\otimes\psi\rangle=0
\]
对充分大的 $M$（依赖于 $|\psi\rangle$）成立。存在正整数 $m_2$，使得当
$k\ge m_2$ 时 $X_\theta(k)|\phi_1\rangle=0$。让我们选取亚纯函数
$h\in H^0(C,\cO_C(Q_{N-1}+*Q_N))$，它在 $Q_{N-1}$ 处具有 Laurent 展开
\[
h_{N-1}=\frac{1}{\xi_{N-1}}+a_{m_2}\xi_{N-1}^{m_2}+\cdots.
\]
于是我们有 $X[f]|\phi_1\rangle=X_\theta(-1)|\phi_1\rangle$。由定理 2.5，存在正整数
$m_3$，使得
\[
X_\theta(-1)^{k}|v_i\rangle=0,\quad i=1,\dots,N-2,\qquad
X_\theta(-1)^{k}|\phi_1\rangle=0,\qquad\text{对一切 }k\ge m_3/N.
\]
于是对 $n\ge m_3$ 我们有
\[
\begin{aligned}
&\langle\tilde\Phi|v_1\otimes\cdots\otimes v_{N-2}
\otimes X_\theta(-1)^{n}\phi_1\otimes\psi\rangle\\
&\qquad=\langle\tilde\Phi|v_1\otimes\cdots\otimes v_{N-2}
\otimes X_\theta[f]^{n}\phi_1\otimes\psi\rangle\\
&\qquad=(-1)^{n}\sum_{n_1+\cdots+n_N=n}
\frac{n!}{n_1!n_2!\cdots n_N!}\langle\tilde\Psi|
\prod_{j\ne N-2}\rho_j(X_\theta[f])^{n_j}u\otimes0\rangle\\
&\qquad=0.
\end{aligned}
\]
现在令
\[
E=X_{-\theta}(1),\qquad F=X_\theta(-1),\qquad H=[E,F].
\]
则 $\{E,F,H\}$ 张成 $\hat{\mathfrak{g}}$ 的一个李子代数 $\mathfrak{k}$，同构于
$\mathfrak{sl}(2,\C)$。注意我们有
\[
H=[X_{-\theta},X_\theta]\otimes1+(X_{-\theta},X_\theta)c=-H_\theta\otimes1+c.
\]

% [Tutor Note: 原文此处「We may assume that $b_{p+1}\ne0$」与上一页展开式的 $b_{-p-1}$ 不一致（同一段后文又用 $b_{-p-1}$），应为 $b_{-p-1}\ne0$。照原文排印，见报告 ERRATA。]
% [Tutor Note: 从 $V_{\lambda_1}\otimes\cdots\otimes\cH_{\lambda_N}$ 下降到 $\cH_{\lambda_{N-1}}$ 需要验证的正是 $\mathfrak{sl}(2)$ 的最高权理论中的条件：$\cH_{\lambda_{N-1}}$ 由最高权向量 $\lambda_{N-1}$ 与 $X_\theta(-1)$ 的幂生成，而 $X_\theta(-1)^{\ell-(\theta,\lambda_{N-1})+1}$ 作用在最高权向量上为零（$\ell-(\theta,\lambda_{N-1})$ 是该表示的「层数」）。]
% [Tutor Note: 这段证明末尾的显示式在原文中与 (3.31) 几乎逐字相同（出现 $\langle\tilde\Psi|$、$u\otimes0\rangle$、$\prod_{j\ne N-2}$，而按上下文应为 $\langle\tilde\Phi|$、$v_1\otimes\cdots\otimes v_{N-2}\otimes\phi_1\otimes\psi\rangle$、$\prod_{j\ne N-1}$）。这是原文的排印混乱，照排并登记；理解时按上下文改正。]

% [原文 p.60]
\figomit{3.2}{正规化。}

因此我们有
\[
H|\lambda_{N-1}\rangle=(\ell-(\theta,\lambda_{N-1}))|\lambda_{N-1}\rangle.
\]
设 $U$ 是 Verma 模 $\cM_{\lambda_{N-1}}$ 的由
\[
U:=\sum_{k=1}^{\infty}\C X_\theta(-1)^{k}|\lambda_{N-1}\rangle
\]
定义的向量子空间。

容易证明 $U$ 是一个 $\mathfrak{k}$-模。设 $U_0$ 是 $U$ 的由
\[
U_0=\{|u\rangle\in U\mid
\langle\tilde\Phi|v_1\otimes\cdots\otimes v_{N-2}\otimes u\otimes\psi\rangle=0\}
\]
定义的子空间，其中 $v_i\in V_i$、$i=1,\dots,N-2$，$|\psi\rangle\in\cH_{\lambda_N}$。则 $U_0$ 是 $\mathfrak{k}$-子模。由于对充分大的 $m\ge m_3$ 我们有 $F^{m}|0\rangle\in U_0$，商空间 $U/U_0$ 是有限维 $\mathfrak{k}$-模。设 $|\overline{0}\rangle$ 是 $|0\rangle\in U_0$ 在 $U/U_0$ 中的像。于是我们有
\[
H|\overline{\lambda_{N-1}}\rangle
=(\ell-(\theta,\lambda_{N-1}))|\overline{\lambda_{N-1}}\rangle.
\]
因此，由 $\mathfrak{sl}(2,\C)$ 的表示论我们有
\[
F^{\ell-(\theta,\lambda)+1}|\overline{\lambda_{N-1}}\rangle=0.
\]
这意味着
\[
\langle\tilde\Psi|u\otimes X_\theta(-1)^{\ell-(\theta,\lambda_{N-1})+1}
|\lambda_{N-1}\rangle=0.
\]
于是我们得到
$\langle\tilde\Psi|\in\Hom(V_1\otimes\cdots\otimes V_{\lambda_{N-2}}
\otimes\cH_{\lambda_{N-1}}\otimes\cH_{\lambda_N},\C)$。

重复这一过程，我们可以唯一地构造出
$\langle\tilde\Phi|\in\cV^{\dagger}_{\vec{\lambda}}$。\hfill$\square$

\noindent\textbf{3.3.2 因子化.}\quad 现在让我们考虑因子化。对元素
$\mu\in P_\ell$，置
\[
\mu^{\dagger}=-w(\mu),
\]
其中 $w$ 是单李代数 $\mathfrak{g}$ 的 Weyl 群的最长元素（换言之，$w(\Delta_+)=\Delta_-$）。注意 $\mu^{\dagger}$ 也由如下事实刻画：$-\mu^{\dagger}$ 是 $\mathfrak{g}$-模 $V_\mu$ 的最低权。例如，对 $\mathfrak{g}=\mathfrak{sl}(2,\C)$，我们有 $\mu^{\dagger}=\mu$，

% [Tutor Note: 原文「$\langle\tilde\Psi|u\otimes X_\theta(-1)^{\ell-(\theta,\lambda_{N-1})+1}|\lambda_{N-1}\rangle=0$」与「$\langle\tilde\Psi|\in\Hom(V_1\otimes\cdots\otimes V_{\lambda_{N-2}}\otimes\cH_{\lambda_{N-1}}\otimes\cH_{\lambda_N},\C)$」中的 $\langle\tilde\Psi|$、$u$ 与 $V_{\lambda_{N-2}}$ 均与上下文（应为 $\langle\tilde\Phi|$、$v_1\otimes\cdots\otimes v_{N-2}$、$V_{\lambda_{N-2}}$ 的记号约定 $V_i=V_{\lambda_i}$）不符，属原文排印混用；照排并登记。]
% [Tutor Note: $U_0$ 的定义里 $\langle\tilde\Phi|$ 作用于 $v_1\otimes\cdots\otimes v_{N-2}\otimes u\otimes\psi$，其中 $u\in U\subset\cM_{\lambda_{N-1}}$；这正是把「对 $\cM_{\lambda_{N-1}}$ 的依赖」压到 $\mathfrak{sl}(2)$ 子模 $U$ 上的关键一步。]
% [Tutor Note: 3.3.2 起进入\textbf{因子化}：这里的 $\mu^{\dagger}=-w(\mu)$ 是「反最高权」操作，$-\mu^\dagger$ 是 $V_\mu$ 的最低权。它是把 $N$ 点共形块拆成 $N-1$ 点共形块与三点函数（OPE 系数）的张量积的起点，也是第 5 章因子化定理的第一次出现。]

% [原文 p.61]
\figomit{3.3}{正规化之后，在曲面内部运动的粒子从 $P_-$ 出来，并在 $P_+$ 处进入曲线。}

\begin{theorem}[3.19]
（\emph{因子化}）在上述记号下，对带形式邻域的稳定 $N$ 点曲线
\[
\mathfrak{X}=(C;Q_1,\dots,Q_N;\eta_1,\dots,\eta_N),
\]
置
\[
\tilde{\mathfrak{X}}=(\tilde{C};P',P'',Q_1,\dots,Q_N;
\eta',\eta'',\eta_1,\dots,\eta_N).
\]
则存在典范同构
\[
\bigoplus_{\mu\in P_\ell}\cV^{\dagger}_{\mu,\mu^{\dagger},\vec{\lambda}}
(\tilde{\mathfrak{X}})\cong\cV^{\dagger}_{\vec{\lambda}}(\mathfrak{X}).
\]
\end{theorem}

\noindent\textit{证明.}\quad $\mathfrak{g}$ 在 $V_\mu\otimes V_{\mu^{\dagger}}$ 上的对角作用使
$V_\mu\otimes V_{\mu^{\dagger}}$ 成为一个 $\mathfrak{g}$-模，并且它包含平凡
$\mathfrak{g}$-模，重数为 $1$。设 $|0_{\mu,\mu^{\dagger}}\rangle$ 是
$V_\mu\otimes V_{\mu^{\dagger}}$ 的平凡 $\mathfrak{g}$-子模的基。也就是说，
\[
\rho_1(X)|0_{\mu,\mu^{\dagger}}\rangle
+\rho_2(X)|0_{\mu,\mu^{\dagger}}\rangle=0
\]
对一切 $X\in\mathfrak{g}$ 成立。

例如，对 $\mathfrak{g}=\mathfrak{sl}(2,\C)$ 与 $\mu=m\alpha$，我们令
$e_0$ 为最高权向量，$He_0=\frac{m}{2}e_0$。置
\[
e_j=\frac{1}{j!}F^{j}e_0,\qquad j=1,2,\dots,m.
\]
于是我们可以选取
\[
|0_{\mu,\mu^{\dagger}}\rangle:=\sum_{j=0}^{m}(-1)^{m}
\frac{j!(m-j)!}{m!}e_j\otimes e_{m-j}.
\]
现在模 $\cH_{\mu,\mu^{\dagger},\vec{\lambda}}$ 包含子空间
\[
\cH_{\mu,\mu^{\dagger},\vec{\lambda}}\supset
|0_{\mu,\mu^{\dagger}}\rangle\otimes\cH_{\vec{\lambda}}
\cong\cH_{\vec{\lambda}}.
\]
对任意元素
$\langle\tilde\Psi|\in\cV^{\dagger}_{\mu,\mu^{\dagger},\vec{\lambda}}
(\tilde{\mathfrak{X}})$，通过
\[
\langle\Psi|\Phi\rangle
=\langle\tilde\Psi|0_{\mu,\mu^{\dagger}}\otimes\Phi\rangle
\]
定义 $\langle\Psi|\in\cH^{\dagger}_{\vec{\lambda}}$。

% [Tutor Note: 这一节是全书第一次出现\textbf{因子化定理}：$\tilde{\mathfrak X}$ 是把原来的 $N$ 点曲线在 $P',P''$ 处「剪开」再接上一点（图 3.3 的正规化），共形块被分解为对所有 $\mu\in P_\ell$ 的直和，每一项对应中间态 $\mu$ 与它的反最高权 $\mu^{\dagger}$。]
% [Tutor Note: 原文单态（singlet）公式 $|0_{\mu,\mu^{\dagger}}\rangle=\sum_{j=0}^{m}(-1)^{m}\frac{j!(m-j)!}{m!}e_j\otimes e_{m-j}$ 中的 $(-1)^{m}$ 应为 $(-1)^{j}$：逐项用 $E$ 作用并令 $e_{j-1}\otimes e_{m-j}$ 的系数为零，得到递推 $c_j=-c_{j-1}j/(m-j+1)$，其解为 $c_j=(-1)^{j}j!(m-j)!/m!$。若取 $(-1)^m$，则 $(-1)^m$ 是整体常数，所得向量不是不变量。照原文排印，登记为 ERRATA。]

% [原文 p.62]
对一切 $|\Phi\rangle\in\cH_{\vec{\lambda}}$。于是，对任意亚纯函数
$f\in H^0(C,\cO_C(*\sum_{j=1}^{N}Q_j))$ 我们有
\[
\begin{aligned}
\sum_{j=1}^{N}\langle\Psi|\rho_j(X[f])|\Phi\rangle
&=\sum_{j=1}^{N}\langle\tilde\Psi|
(0_{\mu,\mu^{\dagger}})\otimes\rho_j(X[f])\Phi\rangle\\
&=\sum_{j=1}^{N+2}\langle\tilde\Psi|
\rho_j(X[f])|0_{\mu,\mu^{\dagger}}\otimes\Phi\rangle=0
\end{aligned}
\]
\begin{sloppypar}这是因为，若把 $f$ 视为 $\tilde{C}$ 上的亚纯函数，则
$f(P')=f(P'')$，并且
$\rho_{P'}(X[f])|0_{\mu,\mu^{\dagger}}\rangle
+\rho_{P''}(X[f])|0_{\mu,\mu^{\dagger}}\rangle=0$。因此我们有
\end{sloppypar}
\[
\sum_{j=1}^{N}\langle\Psi|\rho_j(X[f])=0
\qquad\text{对任意 }f\in H^0(C,\cO_C(*\sum_{j=1}^{N}Q_j)).
\]
于是我们得到典范 $\C$-线性映射
\[
\iota_{\mu}:\cV_{\mu,\mu^{\dagger},\vec{\lambda}}(\tilde{\mathfrak{X}})
\longrightarrow\cV^{\dagger}_{\vec{\lambda}}(\mathfrak{X}).
\]
让我们证明映射 $\iota_{\mu}$ 是单射。为此，首先我们证明对
$\langle\Psi|=\iota_{\mu}(\langle\tilde\Psi|)$，
$\langle\tilde\Psi|\in\cV_{\mu,\mu^{\dagger},\vec{\lambda}}
(\tilde{\mathfrak{X}})$，我们有
\begin{equation}
\langle\Psi|X(P)|\Phi\rangle\dd P
=\langle\tilde\Psi|X(P)|0_{\mu,\mu^{\dagger}}\otimes\Phi\rangle\dd P. \tag{3.33}
\end{equation}
注意由定理 3.15 证明中的断言 3，左端在 $Q_j$ 处关于形式参数
$\xi_j$ 的展开具有形式
\[
\sum_{n\in\Z}\langle\Psi|\rho_j(X(n))|\Phi\rangle
\xi_j^{-n-1}\dd\xi_j.
\]
类似地，(3.33) 的右端具有展开
\[
\begin{aligned}
&\sum_{n\in\Z}\langle\tilde\Psi|\rho_j(X(n))
|0_{\mu,\mu^{\dagger}}\otimes\Phi\rangle\xi_j^{-n-1}\dd\xi_j\\
&\qquad=\sum_{n\in\Z}\langle\tilde\Psi|(0_{\mu,\mu^{\dagger}})
\otimes\rho_j(X(n))\Phi\rangle\xi_j^{-n-1}\dd\xi_j\\
&\qquad=\sum_{n\in\Z}\langle\Psi|\rho_j(X(n))|\Phi\rangle
\xi_j^{-n-1}\dd\xi_j.
\end{aligned}
\]
因此等式 (3.33) 成立。类似的论证给出等式
\[
\begin{aligned}
&\langle\Psi|X_1(P_1)\cdots X_M(P_M)|\Phi\rangle
\dd P_1\cdots\dd P_M\\
&\qquad=\langle\tilde\Psi|X_1(P_1)\cdots X_M(P_M)
|0_{\mu,\mu^{\dagger}}\otimes\Phi\rangle\dd P_1\cdots\dd P_M.
\end{aligned}
\]
现在假设 $\langle\Psi|=0$。由定理 3.15 的第 3) 条我们有
\[
\langle\tilde\Psi|X_2(P_2)\cdots X_M(P_M)
|\rho_{P'}(X_1(n))0_{\mu,\mu^{\dagger}}\otimes\Phi\rangle=0.
\]
再次应用定理 3.15 的第 3) 条，我们得到
\[
\begin{aligned}
\langle\tilde\Psi|\rho_{P'}(X_2(n_2)X_1(n_1))
0_{\mu,\mu^{\dagger}}\otimes\Phi\rangle&=0\\
\langle\tilde\Psi|\rho_{P'}(X_1(n_1))
\rho_{P''}(X_2(n_2))0_{\mu,\mu^{\dagger}}\otimes\Phi\rangle&=0\\
\langle\tilde\Psi|\rho_{P''}(X_1(n_1)X_2(n_2))
0_{\mu,\mu^{\dagger}}\otimes\Phi\rangle&=0.
\end{aligned}
\]
重复同样的过程我们可以证明
\[
\langle\tilde\Psi|\tilde\Phi\rangle=0
\qquad\text{对任意 }\tilde\Phi\in\cH_{\mu,\mu^{\dagger},\vec{\lambda}}
(\tilde{\mathfrak{X}}).
\]

% [Tutor Note: 原文在 $\iota_\mu$ 的定义处及此后数行把定义域印作
%   $\cV_{\mu,\mu^{\dagger},\vec{\lambda}}(\tilde{\mathfrak{X}})$（无上标 $\dagger$），而按 p.60 的记号与上下文应为
%   $\cV^{\dagger}_{\mu,\mu^{\dagger},\vec{\lambda}}(\tilde{\mathfrak{X}})$。照原文排印，登记为 ERRATA 存疑项。]
% [Tutor Note: (3.33) 的证明是「逐点比较展开系数」：左端在 $Q_j$ 处的 Laurent 系数是 $\langle\Psi|\rho_j(X(n))|\Phi\rangle$，右端是 $\langle\tilde\Psi|\rho_j(X(n))|0_{\mu,\mu^{\dagger}}\otimes\Phi\rangle$；而由 $\langle\Psi|$ 的定义，后者等于前者。这一步把「$\langle\Psi|$ 与 $\langle\tilde\Psi|$ 在形式上相等」提升为「作为微分形式逐点相等」，是后面用定理 3.15 第 3) 条做单射性论证的桥梁。]

% [原文 p.63]
因为 $\cH_\mu\otimes\cH_{\mu^{\dagger}}$ 是不可约
$\hat{\mathfrak{g}}\times\hat{\mathfrak{g}}$-模。因此 $\iota_\mu$ 是单射。让我们考虑 $\C$-线性同态
\[
\iota:\bigoplus_{\mu\in P_\ell}
\cV_{\mu,\mu^{\dagger},\vec{\lambda}}(\tilde{\mathfrak{X}})
\xrightarrow{\oplus\iota_\mu}\cV^{\dagger}_{\vec{\lambda}}(\mathfrak{X}).
\]
我们将证明 $\iota$ 是单射。为此，我们对点 $P'$ 与 $P''$ 联系某些右
$\mathfrak{g}$-模与可积右 $\hat{\mathfrak{g}}$-模。

固定元素 $\langle\Psi|\in\cV^{\dagger}_{\vec{\lambda}}(\mathfrak{X})$。设
$h\in H^0(\tilde{C},\cO_{\tilde{C}}(*\sum_{j=1}^{N}Q_j))$ 满足条件
\begin{equation}
h(P')=1,\qquad h(P'')=0. \tag{3.34}
\end{equation}
若 $h'$ 也满足性质 (3.34)，则 $h-h'$ 可以视为 $C$ 上的亚纯函数，且
$h-h'\in H^0(C,\cO_C(*\sum_{j=1}^{N}Q_j))$。因此，由规范条件，对每个
$|u\rangle\in V_{\vec{\lambda}}$，
\[
\sum_{j=1}^{N}\langle\Psi|\rho_j(X[h])|u\rangle
\]
与满足 (3.34) 的亚纯函数 $h$ 的选取无关。对每个元素
$X\in\mathfrak{g}$，通过
\[
\langle\Psi|\rho_{P'}(X)|u\rangle
=-\sum_{j=1}^{N}\langle\Psi|\rho_j(X[h])|u\rangle,
\qquad |u\rangle\in V_{\vec{\lambda}},
\]
定义 $\langle\Psi|\rho_{P'}(X)\in\Hom_{\C}(V_{\vec{\lambda}},\C)$，其中
$h$ 满足 (3.34)。这是良定义的。

其次，对 $X,Y\in\mathfrak{g}$，通过
\[
\langle\Psi|\rho_{P'}(X)\rho_{P'}(Y)|u\rangle
=\sum_{j_1=1,j_2=1}^{N}\langle\Psi|
\rho_{j_2}(Y[h_2])\rho_{j_1}(X[h_1])|u\rangle,
\qquad\text{对一切 }|u\rangle\in V_{\vec{\lambda}},
\]
定义 $\langle\Psi|\rho_{P'}(X)\rho_{P'}(Y)\in\Hom_{\C}
(V_{\vec{\lambda}},\C)$，其中 $h_1$ 与 $h_2$ 满足 (3.34)。由与上面相同的理由，该定义与 $h_2$ 的选取无关。该定义与 $h_1$ 的选取无关可如下证明。由于 $h_2\dd h_1$ 是 $\tilde{C}$ 上的亚纯 $1$-形式，其极点只可能在 $Q_1,\dots,Q_N$ 处，我们有
$\sum_{j=1}^{N}\Res_{Q_j}(h_2\dd h_1)=0$。因此我们有等式
\[
\begin{aligned}
&\sum_{j_1,j_2=1}^{N}\langle\Psi|
\rho_{j_2}(Y[h_2])\rho_{j_1}(X[h_1])|u\rangle\\
&\qquad=\sum_{j_1,j_2=1}^{N}\langle\Psi|
\rho_{j_1}(X[h_1])\rho_{j_2}(Y[h_2])|u\rangle
-\sum_{j=1}^{N}\langle\Psi|\rho_j([X,Y][h_1h_2])|u\rangle.
\end{aligned}
\]
等式的右端表明该定义与 $h_1$ 的选取无关，因为 $h_1h_2$ 也满足性质
(3.34)。此外，上述等式给出等式
\[
\langle\Psi|\bigl(\rho_{P'}(X)\rho_{P'}(Y)-\rho_{P'}(Y)\rho_{P'}(X)\bigr)
=\langle\Psi|\rho_{P'}([X,Y]).
\]
这样我们便可在点 $P'$ 处定义右 $\mathfrak{g}$-模
$U(\langle\Psi|)\subset\Hom_{\C}(V_{\vec{\lambda}},\C)$。用同样的方式，我们可以在点 $P''$ 处构造右 $\mathfrak{g}$-模，其中 $M$ 是充分大的整数。

% [原文 p.64]
我们也可以定义可积右 $\hat{\mathfrak{g}}$-模
$\hat{U}(\langle\Psi|)\subset\Hom_{\C}(\cH_{\vec{\lambda}},\C)$。例如，
$\langle\Psi|\rho_{P'}(X(n))$ 如下定义。设 $g$ 是 $\tilde{C}$ 上的亚纯函数，使得
\begin{equation}
\begin{gathered}
g\in H^0\bigl(\tilde{C},\cO_{\tilde{C}}
(*(\sum_{j=1}^{N}Q_j+P'))\bigr)\\
g\equiv\xi'^{n}\mod(\xi'^{M})\qquad\text{在 }P'\text{ 处}\\
g(P'')=0
\end{gathered}\tag{3.35}
\end{equation}
其中 $\xi'=\eta'^{-1}(\xi)$ 是点 $P'$ 处的形式参数。通过
\[
\langle\Psi|\rho_{P'}(X(n))|u\rangle
=-\sum_{j=1}^{N}\langle\Psi|\rho_j(X[g])|u\rangle
\]
定义 $\langle\Psi|\rho_{P'}(X(n))$。

\begin{sloppypar}该定义与满足 (3.35) 的亚纯函数 $g$ 的选取无关。类似地，我们可以定义
$\langle\Psi|\rho_{P'}(X(n))\rho_{P'}(Y(m))\in
\Hom_{\C}(\cH_{\vec{\lambda}},\C)$，并且我们有等式
\end{sloppypar}
\begin{equation}
\begin{aligned}
&\langle\Psi|\bigl(\rho_{P'}(X(n))\rho_{P'}(Y(m))
-\langle\Psi|\rho_{P'}(Y(m))\rho_{P'}(X(n))\bigr)\\
&\qquad=\langle\Psi|\rho_{P'}([X,Y](m+n))
+\ell\cdot(X,Y)n\delta_{n+m,0}\langle\Psi|.
\end{aligned}\tag{3.36}
\end{equation}
这样我们便可以构造右 $\hat{\mathfrak{g}}$-模
$\hat{U}(\langle\Psi|)\subset\Hom_{\C}(\cH_{\vec{\lambda}},\C)$。
由于由定理 2.5，$\rho_j(X_\alpha[g])$（$X_\alpha\in\mathfrak{g}_\alpha$）的作用是局部幂零的，所以 $\rho_{P'}(X_\alpha(m))$ 在 $\hat{U}(\langle\Psi|)$ 上的作用是局部幂零的。因此 $\hat{U}(\langle\Psi|)$ 是水平为 $\ell$ 的可积右 $\hat{\mathfrak{g}}$-模。

由于 $\cV^{\dagger}_{\vec{\lambda}}(\mathfrak{X})$ 是有限维的，我们可以对点 $P'$ 联系一个有限右 $\mathfrak{g}$-模
\[
U(\cV^{\dagger}_{\vec{\lambda}}(\mathfrak{X}))
=\bigcup_{\langle\Psi|\in\cV^{\dagger}_{\vec{\lambda}}(\mathfrak{X})}
U(\langle\Psi|)
\subset\Hom_{\C}(V_{\vec{\lambda}},\C)
\]
与一个水平为 $\ell$ 的可积右 $\hat{\mathfrak{g}}$-模
\[
\hat{U}(\cV^{\dagger}_{\vec{\lambda}}(\mathfrak{X}))
=\bigcup_{\langle\Psi|\in\cV^{\dagger}_{\vec{\lambda}}(\mathfrak{X})}
\hat{U}(\langle\Psi|)
\subset\Hom_{\C}(\cH_{\vec{\lambda}},\C).
\]
由 (3.36) 我们有不可约分解
\[
U(\cV^{\dagger}_{\vec{\lambda}}(\mathfrak{X}))
=\bigoplus_{\mu\in P_\ell}V_{\mu^{\dagger}}^{\oplus n_\mu}.
\]
于是，由引理 2.15 之前给出的注，我们有
\begin{equation}
\hat{U}(\cV^{\dagger}_{\vec{\lambda}}(\mathfrak{X}))
=\bigoplus_{\mu\in P_\ell}\cH^{\dagger\oplus n_\mu}_{\mu}. \tag{3.37}
\end{equation}
现在我们可以证明 $\iota$ 的单射性了。对元素
$\langle\tilde\Psi|\in\cV_{\mu,\mu^{\dagger},\vec{\lambda}}
(\tilde{\mathfrak{X}})$，置 $\langle\Psi|=\iota_{\mu}(\langle\tilde\Psi|)$，并选取 $\tilde{C}$ 上满足 (3.34) 的亚纯函数 $h$。

% [原文 p.65]
于是我们有
\[
\begin{aligned}
&\langle\Psi|\rho_{P'}(X_1)\cdots\rho_{P'}(X_k)|u\rangle\\
&\qquad=(-1)^{k}\sum_{j_1=1,\cdots,j_k=1}^{N}\langle\Psi|
\rho_{j_1}(X_1[h])\cdots\rho_{j_k}(X_k[h])|u\rangle\\
&\qquad=(-1)^{k+1}\langle\tilde\Psi|\rho_{P'}(X_1(0))
\cdots\rho_{P'}(X_k(0))|0_{\mu,\mu^{\dagger}}\otimes u\rangle.
\end{aligned}
\]
由于
$\rho_{P'}(X_1(0))\cdots\rho_{P'}(X_k(0))|0_{\mu,\mu^{\dagger}}\rangle$
生成同构于 $V_{\mu^{\dagger}}$ 的不可约左 $\mathfrak{g}$-模，我们断言
\[
U(\langle\Psi|)\subset V_{\mu}^{\dagger\oplus n_\mu}.
\]
因此，对
$\langle\tilde\Psi_\mu|\in\cV_{\mu,\mu^{\dagger},\vec{\lambda}}
(\tilde{\mathfrak{X}})$ 与
$\langle\tilde\Psi_\nu|\in\cV_{\nu,\nu^{\dagger},\vec{\lambda}}
(\tilde{\mathfrak{X}})$，$\mu\ne\nu$，我们有
\[
U(\langle\tilde\Psi_\mu|)\cap U(\langle\tilde\Psi_\nu|)=0.
\]
这意味着 $\iota$ 是单射，因为 $\iota_\mu$ 是单射。

最后我们证明 $\iota$ 是满射。由 (3.37)，对元素
$\langle\Psi|\in\cV^{\dagger}_{\vec{\lambda}}(\mathfrak{X})$ 我们有分解
\[
\langle\Psi|=\sum_{\mu\in P_\ell}\langle\Psi_\mu|,
\qquad\langle\Psi_\mu|\in V_{\mu}^{\dagger\oplus n_\mu}.
\]
我们如下构造
$\langle\tilde\Psi_\mu|\in\Hom_{\C}
(\cH_{\mu,\mu^{\dagger},\vec{\lambda}},\C)$。首先注意
$V_\mu\otimes V_{\mu^{\dagger}}$ 由元素
\[
\rho_{P'}(X_1)\cdots\rho_{P'}(X_n)\rho_{P''}(Y_1)\cdots
\rho_{P''}(Y_m)|0_{\mu,\mu^{\dagger}}\rangle,
\qquad X_1,\dots,X_n,Y_1,\dots,Y_m\in\mathfrak{g}
\]
生成。此外，$\langle\tilde\Psi_\mu|$ 定义一个右
$\hat{\mathfrak{g}}$-模
$\hat{U}(\langle\tilde\Psi_\mu|)\subset\Hom_{\C}
(\cH_{\vec{\lambda}},\C)$。对每个元素 $|v\rangle\in\cH_{\vec{\lambda}}$ 定义
\[
\langle\tilde\Psi_\mu|0_{\mu,\mu^{\dagger}}\otimes v\rangle
=\langle\Psi|v\rangle.
\]
定义
\[
\begin{aligned}
&\langle\tilde\Psi_\mu|\rho_{P'}(X_1)\cdots\rho_{P'}(X_n)
\rho_{P''}(Y_1)\cdots\rho_{P''}(Y_m)
|0_{\mu,\mu^{\dagger}}\otimes v\rangle\\
&\qquad=(-1)^{m}\langle\Psi_\mu|\rho_{P'}(X_1(0))\cdots
\rho_{P'}(X_n(0))\rho_{P'}(Y_m(0))\cdots\rho_{P'}(Y_1(0))v\rangle.
\end{aligned}
\]
这是良定义的，因为 $\mathfrak{g}$ 在
$\C|0_{\mu,\mu^{\dagger}}\rangle$ 上的对角作用是平凡的。这就定义了
$\langle\tilde\Psi_\mu|\in\Hom_{\C}
(V_\mu\otimes V_{\mu^{\dagger}}\otimes\cH_{\vec{\lambda}},\C)$。现在假设我们已经定义了
$\langle\tilde\Psi_\mu|\in\Hom_{\C}(F_p\cH_\mu\otimes
F_q\cH_{\mu^{\dagger}}\otimes\cH_{\vec{\lambda}},\C)$，其中 $p,q$ 为非负整数。取元素
$\rho_{P'}(X(m))|u\otimes u'\otimes v\rangle\in
F_{p+1}\cH_\mu\otimes F_q\cH_{\mu^{\dagger}}$，其中
$|u\otimes u'\rangle\in F_p\cH_\mu\otimes F_q\cH_{\mu^{\dagger}}$。选取
$\tilde{C}$ 上的亚纯函数 $f$，使得
\[
\begin{aligned}
f&\in H^0\bigl(\tilde{C},\cO_{\tilde{C}}
(*\sum_{j=1}^{N}Q_j+*P'+*P'')\bigr)\\
f&\equiv\xi'^{m}\mod(\xi'^{M})\qquad\text{在 }P'\text{ 处}\\
f&\equiv0\mod(\xi''^{M})\qquad\text{在 }P''\text{ 处}.
\end{aligned}
\]
这里我们这样选取正整数 $M$，使得
\[
\rho_{P'}(X(n))|u\rangle=0,\qquad
\rho_{P''}(X(n))|u\rangle=0
\]

% [原文 p.66]
对一切 $n\ge M$。然后我们定义
\[
\langle\tilde\Psi_\mu|\rho_{P'}(X(m))|u\otimes u'\otimes v\rangle
=-\sum_{j=1}^{N}\langle\tilde\Psi_\mu|\rho_j(X[f])
|u\otimes u'\otimes v\rangle.
\]
通过与定理 3.18 的证明类似的论证，我们可以证明该定义与满足上述条件的亚纯函数 $f$ 的选取无关，并且我们有
\[
\langle\tilde\Psi_\mu|\in\Hom_{\C}
(F_{p+1}\cH_\mu\otimes F_q\cH_{\mu^{\dagger}}
\otimes\cH_{\vec{\lambda}},\C).
\]
类似地我们可以定义
\[
\langle\tilde\Psi_\mu|\in\Hom_{\C}
(F_p\cH_\mu\otimes F_{q+1}\cH_{\mu^{\dagger}}
\otimes\cH_{\vec{\lambda}},\C).
\]
这样我们便可以证明存在
\[
\langle\tilde\Psi_\mu|\in\Hom_{\C}
(\cH_\mu\otimes\cH_{\mu^{\dagger}}\otimes\cH_{\vec{\lambda}},\C).
\]
此外，我们可以验证规范条件，使得
\[
\langle\tilde\Psi_\mu|\in\cV_{\mu,\mu^{\dagger},\vec{\lambda}}
(\tilde{\mathfrak{X}}).
\]
由我们的构造有
$\iota_\mu(\langle\tilde\Psi_\mu|)=\langle\Psi_\mu|$。\hfill$\square$

\begin{corollary}[3.20]
存在典范同构
\[
\cV_{\vec{\lambda}}(\mathfrak{X})
\cong\bigoplus_{\mu\in P_\ell}
\cV_{\mu,\mu^{\dagger},\vec{\lambda}}(\tilde{\mathfrak{X}}).
\]
\end{corollary}

\begin{remark}[3.21]
在第 6 章我们将证明
$\dim\cV^{\dagger}_{\vec{\lambda}}(\mathfrak{X})$ 只依赖于曲线的亏格 $g$ 与
$\vec{\lambda}$。因此，通过把光滑曲线退化到带结点的曲线并利用上述定理
3.3.5，$\dim\cV^{\dagger}_{\vec{\lambda}}(\mathfrak{X})$ 的计算就归结为亏格 $0$ 曲线的情形。这将在第 6 章讨论。
\end{remark}

对亏格 $0$ 的 $N$ 点光滑曲线，定理 3.18 可以用如下更强的陈述代替。

\begin{proposition}[3.22]
若 $\mathfrak{X}$ 的底曲线是 $\PP^{1}$，则典范映射
\[
\cV^{\dagger}_{\vec{\lambda}}(\mathfrak{X})
\longrightarrow\Hom_{\mathfrak{g}}(V_{\vec{\lambda}},\C)
\]
是单射。
\end{proposition}

证明将在 §6.1 中给出。注意若底曲线的亏格为正，命题 3.22 不成立。例如，对带形式邻域的单点椭圆曲线
\[
\mathfrak{X}=(E;Q,\zeta),
\]
正如我们将在 §5.4 中看到的，我们有
\[
\dim\cV^{\dagger}_{0}(\mathfrak{X})=\ell+1.
\]
但 $\Hom(V_0,\C)$ 是一维的。

作为命题 3.22 的一个弱推广，我们有如下命题。

\begin{proposition}[3.23]
若 $N\ge g+1$，则自然映射
\[
\cV^{\dagger}_{\vec{\lambda}}(\mathfrak{X})
\longrightarrow\Hom_{\mathfrak{g}}
\bigl(F_1(\otimes_{j=1}^{N}\cH_{\lambda_j}),\C\bigr)
\]
是单射。
\end{proposition}

% [Tutor Note: 命题 3.22 与 3.23 的意义：$\cV^\dagger_{\vec\lambda}(\mathfrak X)$ 中的元素若在 $\Hom_{\mathfrak g}(V_{\vec\lambda},\C)$（或 $\Hom_{\mathfrak g}(F_1(\otimes\cH_{\lambda_j}),\C)$）中的像为零，则本身为零。这给出共形块空间维数的上界。命题 3.22 只对 $\PP^1$ 成立；正文用单点椭圆曲线 $\dim\cV^\dagger_0=\ell+1$ 与 $\dim\Hom(V_0,\C)=1$ 的对比说明亏格正时单射失效。]

